\begin{flushleft}
\begin{table}[h!]
\centering
\begin{tabularx}{\textwidth}{ |X|X| }
\hline
Criterium & Voldoende als ... \\
\hline
Theorie	Student & legt helder uit wat een herstelpunt is en kan het verschil met back-up en systeemimage correct benoemen. \\
\hline
Praktijk & Herstelpunt is aangemaakt, correct benoemd en aantoonbaar gebruikt. \\
\hline
Testen	& Het terugzetten van een herstelpunt is correct uitgevoerd en getest. \\
\hline
Rapportage & Bevat duidelijke uitleg, screenshots en een reflectie op voordelen/nadelen. \\
\hline
Professionaliteit & Student werkt netjes, verantwoord en kan vragen van docent beantwoorden. \\

\hline
\end{tabularx}
\caption{Beoordelingscriteria}
\label{table:facility}
\end{table}
\end{flushleft}

Extra (voor verdieping)
\begin{itemize}
\item Onderzoek hoe herstelpunten via PowerShell beheerd kunnen worden.
\item Beschrijf situaties waarin herstelpunten niet voldoende zijn en een echte back-up noodzakelijk is.
\end{itemize}

